%% CPP 2027 Submission — 0pat Exercise XXXVI: Divergence–Period Symmetry on the Zeta Series
\documentclass[sigplan,10pt,anonymous,review]{acmart}
\settopmatter{printfolios=true,printacmref=false}

\AtBeginDocument{%
  \providecommand\BibTeX{{\normalfont B\scshape ib\TeX}}}

\settopmatter{printacmref=false}
\renewcommand\footnotetextcopyrightpermission[1]{}
\pagestyle{plain}

\usepackage{microtype}
\usepackage{booktabs}
\usepackage{amsmath}
% XeTeX (tectonic): acmart loads newtxmath/libertine which already defines \Bbbk;
% reset it so amssymb's definition is accepted.
\let\Bbbk\relax
\usepackage{amssymb}

\begin{document}

\title{The Divergence--Period Symmetry Realized on the Zeta Series:\\
One Conjugation, Two Eigenspaces (Exercise XXXVI)}

\author{Anonymous}
\affiliation{%
  \institution{Anonymous}
}
\email{anonymous@example.org}

\begin{abstract}
We formalize, with five machine-checked theorems, the way in which a
single conjugation symmetry splits the terms of the zeta series into
two eigenspaces. Each term $1/n^s$ factors as $n^{-\sigma} \cdot
e^{-it\ln n}$, where $s = \sigma + it$: the divergence axis (the real
part) controls the modulus, $\|n^s\| = n^{\sigma}$, and the period axis
(the imaginary part) contributes a unit modulus that never affects
convergence. The two axes are the two eigenspaces of the same
symmetry; each axis has its own mirror pairs that collapse to the unit
$1$: $e^{i\theta} \cdot e^{-i\theta} = 1$ on the period axis, and
$r \cdot (1/r) = 1$ on the divergence axis. Conjugation is the reversal
of the period axis, $\overline{n^s} = n^{\overline{s}}$ for positive
real $n$. The functional equation
$\zeta(1-s) = 2(2\pi)^{-s}\Gamma(s)\cos(\pi s/2)\zeta(s)$ pairs the
divergent region with the convergent region through this same symmetry
--- it is the mechanism of analytic continuation, and it is already a
machine-checked theorem of mathlib
(\texttt{riemannZeta\_one\_sub}). Zero errors, zero \texttt{sorry}.
\end{abstract}

\keywords{Lean, formalization, zeta function, functional equation, analytic continuation, conjugation symmetry, 0pat}

\maketitle

\section{The series term as a product of two eigenspaces}

The zeta function is born as a series, and each term of the series is a
product of two factors. For a positive integer $n$ and a complex
exponent $s = \sigma + it$,

\begin{equation*}
  n^s = n^{\sigma + it} = n^{\sigma} \cdot n^{it}
      = n^{\sigma} \cdot e^{it \ln n}.
\end{equation*}

The first factor $n^{\sigma}$ is real; it grows or decays with the real
part of $s$. The second factor $e^{it \ln n}$ lies on the unit circle;
it oscillates with the imaginary part of $s$. The series term
$1/n^s = n^{-\sigma} \cdot e^{-it \ln n}$ is therefore the product of
a \emph{divergence factor} (real, monotone in $\sigma$) and a
\emph{period factor} (unit modulus, oscillating in $t$).

The first theorem makes this split precise and machine-checked.

\paragraph{The modulus of a term is the divergence factor
(\texttt{zeta\_term\_norm\_split}).} For $n > 0$ and $s \in
\mathbb{C}$,

\begin{equation*}
  \|n^s\| = n^{\mathrm{Re}\,s}.
\end{equation*}

The modulus of the term is entirely controlled by the divergence axis
(the real part of $s$). The period axis never affects the modulus:
$|e^{-it\ln n}| = 1$ for every $t$. In this sense the two axes play
asymmetric roles in convergence --- but they are not independent
structures; the next section shows that both are eigenspaces of one
and the same conjugation symmetry.

\section{One symmetry, two mirror pairs}

Conjugation $z \mapsto \overline{z}$ fixes the real axis pointwise and
reverses the imaginary axis. The two axes of the series term are the
two eigenspaces of this single symmetry: a real number is fixed
(eigenvalue $+1$), a pure imaginary number is negated (eigenvalue
$-1$). Each eigenspace carries its own mirror pairs, and on both axes
a mirror pair collapses to the unit $1$.

\paragraph{The period-axis mirror pair (\texttt{period\_pair\_reduces}).}
The pair $e^{i\theta}$ and $e^{-i\theta}$ are conjugate: for real
$\theta$, $\overline{e^{i\theta}} = e^{-i\theta}$. Their product is the
unit,

\begin{equation*}
  e^{i\theta} \cdot e^{-i\theta} = e^{0} = 1 .
\end{equation*}

The exponent coordinate is negated, and the two mirror terms annihilate
to the neutral element of multiplication. This is the algebraic
identity $\exp(\theta) \cdot \exp(-\theta) = 1$, machine-checked for
complex $\theta$.

\paragraph{The divergence-axis mirror pair
(\texttt{divergence\_pair\_reduces}).} The pair $r$ and $1/r$ is not a
conjugate pair in the analytic sense --- both numbers are real --- but
it is the mirror pair of the divergence axis: under the logarithm the
two coordinates are negatives of each other,
$\ln(1/r) = -\ln r$. Their product is the unit,

\begin{equation*}
  r \cdot \frac{1}{r} = 1 \qquad (r \neq 0).
\end{equation*}

The logarithmic coordinate is negated, and the two mirror terms
annihilate to the neutral element of multiplication, exactly as on the
period axis.

\paragraph{The two pairs are the two faces of one symmetry.} On the
period axis the mirror relation is $\theta \mapsto -\theta$ (via
conjugation); on the divergence axis it is $\ln r \mapsto -\ln r$ (via
reciprocation). In both cases a single coordinate is negated and the
product of the mirror pair is $1$. The two theorems
\texttt{period\_pair\_reduces} and \texttt{divergence\_pair\_reduces}
verify both faces: the same collapse, on the two eigenspaces of the
same symmetry.

\section{Conjugation is the reversal of the period axis}

If the series term $n^s$ is a product of a divergence factor and a
period factor, then conjugation acts on it through the period factor
alone. Conjugating the whole term is the same as conjugating the
exponent:

\paragraph{Conjugation passes through the power
(\texttt{term\_conj\_symmetry}).} For $n > 0$ and $s \in \mathbb{C}$,

\begin{equation*}
  \overline{n^s} = n^{\overline{s}} .
\end{equation*}

Since $n$ is positive real, its argument is $0$, not $\pi$; mathlib's
\texttt{Complex.conj\_cpow} then moves the conjugation across the
power. The statement is exactly the eigenspace behavior of the
symmetry: $\overline{s} = \sigma - it$ fixes the real part (the
divergence axis, $\sigma$ unchanged) and negates the imaginary part
(the period axis, $t \mapsto -t$). Conjugation reverses the period
axis and leaves the divergence axis untouched.

\section{The same symmetry pairs the two regions}

The split of the term into two eigenspaces is not only a local
observation about a single term. The same symmetry connects the region
where the series diverges with the region where it converges.

\paragraph{The functional equation is the pairing
(\texttt{functional\_equation\_bridges}).} For $s \neq 1$ and
$s \neq -n$ for all $n \in \mathbb{N}$,

\begin{equation*}
  \zeta(1-s) = 2\,(2\pi)^{-s}\,\Gamma(s)\,
  \cos\!\left(\frac{\pi s}{2}\right)\, \zeta(s) .
\end{equation*}

This is the functional equation of the zeta function, in the exact
form of mathlib's \texttt{riemannZeta\_one\_sub}; the theorem above is
that statement, machine-checked. The map $s \mapsto 1-s$ sends the
convergent region ($\mathrm{Re}\,s > 1$) to the divergent region
($\mathrm{Re}\,s < 0$): the values of the divergent series are obtained
from the convergent region through this symmetry. The pairing $s
\leftrightarrow 1-s$ is the same conjugation symmetry of the series
terms acting on the exponent: $\sigma \mapsto 1-\sigma$ negates the
divergence coordinate around its fixed point, exactly as
$t \mapsto -t$ negates the period coordinate. The mechanism by which
the divergent series acquires its values is the mechanism by which the
two axes were mirrored in the first place.

\section{Honest boundary and position}

The content is classical: the factorization of $n^s$ into a real power
and a unit-modulus phase, the two mirror identities
$e^{i\theta}e^{-i\theta} = 1$ and $r(1/r) = 1$, conjugation across the
power for positive real bases, and the zeta functional equation. All
five theorems are known mathematics, and the fifth is quoted directly
from mathlib rather than proved anew. The value of the exercise is the
precise statement of the structure: the divergence axis and the period
axis of the series term are the two eigenspaces of one conjugation
symmetry, and the functional equation is that same symmetry pairing
the divergent region with the convergent region --- analytic
continuation as the symmetry's mechanism. No claim is made beyond the
five machine-checked statements; the reader is invited to inspect the
artifact.

\section*{Artifact}

The artifact contains \texttt{proof.lean} (Lean 4.32.2, mathlib
v4.32.2): five theorems, zero errors, zero \texttt{sorry}. Build:
\texttt{lean proof.lean}.

\end{document}
